\section{Funktory: odwzorowania między kategoriami}\label{sec:funktory}

\subsection{Definicja i przykłady}

\begin{definicja}[Funktor kowariantny]\label{def:funktor}
Niech $\mathcal C,\mathcal D$ będą kategoriami. \emph{Funktor (kowariantny)} $F:\mathcal C\to\mathcal D$ przyporządkowuje każdemu obiektowi $A$ obiekt $F(A)$ i każdemu morfizmowi $f:A\to B$ morfizm $F(f):F(A)\to F(B)$ w taki sposób, że dla wszystkich odpowiednich morfizmów zachodzi $F(\id_A)=\id_{F(A)}$ oraz $F(g\circ f)=F(g)\circ F(f)$.
\end{definicja}

Funktor zachowuje \emph{sposób organizacji strzałek}. Nie musi zachowywać dosłownie wszystkich dodatkowych struktur, które możemy rozpoznać w wyjściowej kategorii.

\begin{przyklad}[Funktor zapominania]
Funktor $U:\Grp\to\Set$ przypisuje grupie $G$ jej zbiór elementów $U(G)$, a homomorfizmowi $\varphi:G\to H$ tę samą funkcję, traktowaną wyłącznie jako odwzorowanie zbiorów. Niczego nie zmienia w funkcji $\varphi$; jedynie \emph{zapomina}, że funkcja zachowuje mnożenie. Analogicznie mamy $U:\Vect_k\to\Set$ i $U:\Top\to\Set$.
\end{przyklad}

\begin{przyklad}[Fundamentalny przykład z algebry liniowej]
Funktor $F:\cat{Mat}_k\to\Vect_k^{\mathrm{fd}}$ definiujemy przez $F(n)=k^n$ i $F(A)(x)=Ax$. Zachowanie składania jest równością $(BA)x=B(Ax)$, a zachowanie identyczności wynika z $I_nx=x$. Każda skończenie wymiarowa przestrzeń wektorowa jest izomorficzna z pewnym $F(n)$, lecz nie jest dosłownie równa takiej przestrzeni bez wyboru bazy.
\end{przyklad}

\begin{przyklad}[Potęga symetryczna jako funktor]
Ustalmy $r\geq1$. Dla przestrzeni $V$ nad $k$ oznaczmy przez $\operatorname{Sym}^r(V)$ $r$-tą potęgę symetryczną, a dla $f:V\to W$ przez $\operatorname{Sym}^r(f)$ jedyne odwzorowanie liniowe określone na iloczynach symetrycznych wzorem $v_1\cdots v_r\mapsto f(v_1)\cdots f(v_r)$. Odwzorowanie to istnieje, ponieważ $r$-liniowe odwzorowanie $(v_1,\ldots,v_r)\mapsto f(v_1)\cdots f(v_r)$ jest symetryczne, więc faktoryzuje się przez potęgę symetryczną. Iloczyny symetryczne generują tę przestrzeń, a na każdym takim generatorze zarówno $\operatorname{Sym}^r(g\circ f)$, jak i $\operatorname{Sym}^r(g)\circ\operatorname{Sym}^r(f)$ dają $g(f(v_1))\cdots g(f(v_r))$. Dla identyczności każdy generator pozostaje niezmieniony. Zatem $\operatorname{Sym}^r:\Vect_k\to\Vect_k$ jest funktorem. To przykład konstrukcji algebraicznej działającej również na odwzorowaniach.
\end{przyklad}

\begin{przyklad}[Grupa działająca na zbiorze]\label{przyklad:dzialanie-funktor}
Traktujmy grupę $G$ jako kategorię z jednym obiektem $*$, jak w przykładzie~\ref{przyklad:monoid-kat}. Funktor $F:G\to\Set$ wybiera zbiór $X=F(*)$ i dla każdego $g\in G$ bijekcję $F(g):X\to X$ (bijekcję, ponieważ $F(g^{-1})$ jest odwrotnością $F(g)$). Równość $F(gh)=F(g)\circ F(h)$ oznacza dokładnie aksjomat lewego działania $g\cdot(h\cdot x)=(gh)\cdot x$. A zatem \emph{lewe działania grupy $G$ na zbiorach są tym samym co funktory $G\to\Set$}.
\end{przyklad}

\begin{przyklad}[Funktory między kategoriami z jednym obiektem]
Monoidy $M,N$ rozumiemy jako kategorie z jednym obiektem. Funktor $F:M\to N$ jest dokładnie homomorfizmem monoidów: funkcja $M\to N$ musi zachować jednostkę i iloczyn. Jeśli oba monoidy są grupami, otrzymujemy zwykłe homomorfizmy grup. Tak jeden i ten sam aksjomat funktora przybiera postać znajomego warunku algebraicznego.
\end{przyklad}


\begin{przyklad}[Zbiór składowych spójnych jako funktor]
Niech $c(X)$ oznacza zbiór składowych spójnych przestrzeni topologicznej $X$. Oznaczenie $\pi_0(X)$ rezerwujemy dla zbioru składowych łukowej spójności (klas punktów połączonych drogą), które nie zawsze pokrywają się ze składowymi spójnymi. Dla ciągłego $f:X\to Y$ obraz każdej składowej spójnej $K\subseteq X$ jest spójny, więc zawiera się w dokładnie jednej składowej $L\subseteq Y$. Możemy zatem zdefiniować $c(f)(K)=L$. Identyczność pozostawia każdą składową na miejscu. Jeżeli $g:Y\to Z$ jest ciągłe, a $f(K)\subseteq L=c(f)(K)$, to $g(f(K))\subseteq g(L)\subseteq c(g)(L)$, zatem jednoznaczność składowej zawierającej niepusty zbiór spójny daje $c(g\circ f)(K)=c(g)(c(f)(K))$. Otrzymujemy funktor $c:\Top\to\Set$. Dla sumy topologicznej dwóch niepustych przedziałów składowe są dwie; dla przedziału $[0,1]$ jest jedna. Funktor wydobywa niezmiennik topologiczny i pamięta działanie ciągłych odwzorowań na tym niezmienniku.
\end{przyklad}

\begin{przyklad}[Funktor określony przez diagram]
Niech $\mathcal J$ będzie kategorią z obiektami $0,1$ i jedną nietrywialną strzałką $a:0\to1$. Funktor $F:\mathcal J\to\Set$ jest dokładnie wyborem dwóch zbiorów $X=F(0)$, $Y=F(1)$ i funkcji $f=F(a):X\to Y$. Dla $X=\{1,2,3\}$, $Y=\{a,b\}$ i $f(1)=f(2)=a$, $f(3)=b$ otrzymujemy konkretny diagram, który sam jest funktorem. Zmiana kształtu kategorii indeksującej prowadzi do innych typów diagramów.
\end{przyklad}


\subsection{Funktory kontrawariantne}

\begin{definicja}[Funktor kontrawariantny]
\emph{Funktor kontrawariantny} z $\mathcal C$ do $\mathcal D$ jest funktorem kowariantnym $F:\mathcal C^{\op}\to\mathcal D$. W zapisie z morfizmami kategorii $\mathcal C$ oznacza to, że strzałce $f:A\to B$ przypisujemy strzałkę $F(f):F(B)\to F(A)$, przy czym $F(\id_A)=\id_{F(A)}$ oraz $F(g\circ f)=F(f)\circ F(g)$.
\end{definicja}

\begin{przyklad}[Przeciwobraz]
Dla zbioru $X$ niech $\mathcal P(X)$ oznacza zbiór jego podzbiorów. Funkcji $f:X\to Y$ przyporządkujmy $\mathcal P(f):\mathcal P(Y)\to\mathcal P(X)$, $B\mapsto f^{-1}(B)$. Równość $(g\circ f)^{-1}(C)=f^{-1}(g^{-1}(C))$ daje funktor kontrawariantny $\mathcal P:\Set^{\op}\to\Set$.

Istnieje też \emph{inny}, kowariantny funktor potęgowy $\mathcal P_+:\Set\to\Set$, który funkcji $f:X\to Y$ przypisuje $A\mapsto f(A)$. Warto nie mylić tych dwóch konstrukcji.
\end{przyklad}

\begin{przyklad}[Przestrzeń dualna]\label{przyklad:dualny}
Dla przestrzeni wektorowej $V$ niech $V^*=\Hom_k(V,k)$. Odwzorowaniu liniowemu $f:V\to W$ przypisujemy $f^*:W^*\to V^*$, $\lambda\mapsto\lambda\circ f$. Mamy $(g\circ f)^*=f^*\circ g^*$, a zatem $(-)^*:\Vect_k^{\op}\to\Vect_k$ jest funktorem kontrawariantnym. Dla skończenie wymiarowych $V$ jest on szczególnie dobrze związany z podwójną dualnością, którą omówimy w rozdziale~\ref{sec:naturalne}.
\end{przyklad}

\begin{stwierdzenie}\label{tw:funktory-iso}
Każdy funktor $F:\mathcal C\to\mathcal D$ przenosi izomorfizmy na izomorfizmy. Jeśli $f:A\to B$ jest izomorfizmem w $\mathcal C$, to $F(f)^{-1}=F(f^{-1})$.
\end{stwierdzenie}
\begin{proof}
Z $f^{-1}\circ f=\id_A$ otrzymujemy $F(f^{-1})\circ F(f)=F(\id_A)=\id_{F(A)}$. Podobnie $F(f)\circ F(f^{-1})=\id_{F(B)}$.
\end{proof}

\begin{uwaga}
Funktor nie musi przenosić dowolnych monomorfizmów na monomorfizmy ani dowolnych epimorfizmów na epimorfizmy. Powód jest subtelny: w definicji mono- lub epimorfizmu kwantyfikujemy po \emph{wszystkich} morfizmach w danej kategorii, a po przejściu do drugiej kategorii mogą pojawić się nowe morfizmy, których nie widzieliśmy w wyjściowej.
\end{uwaga}

\subsection{Wierność i pełność}

\begin{definicja}[Funktor wierny i pełny]\label{def:pelny}
Funktor $F:\mathcal C\to\mathcal D$ indukuje dla każdego $A,B\in\Ob(\mathcal C)$ funkcję $F_{A,B}:\Hom_{\mathcal C}(A,B)\to\Hom_{\mathcal D}(F(A),F(B))$, $f\mapsto F(f)$. Funktor jest \emph{wierny}, jeśli wszystkie $F_{A,B}$ są różnowartościowe; jest \emph{pełny}, jeśli wszystkie $F_{A,B}$ są surjektywne; jest \emph{w pełni wierny}, jeśli wszystkie są bijekcjami.
\end{definicja}

\begin{przyklad}
Funktor zapominania $U:\Grp\to\Set$ jest wierny: dwa homomorfizmy o tej samej funkcji na elementach są równe. Nie jest pełny, gdyż nie każda funkcja między zbiorami bazowymi grup jest homomorfizmem.
\end{przyklad}

\begin{przyklad}
Włączenie pełnej podkategorii $\Vect_k^{\mathrm{fd}}\hookrightarrow\Vect_k$ jest w pełni wierne. Funktor $\cat{Mat}_k\to\Vect_k^{\mathrm{fd}}$ z poprzedniego przykładu również jest w pełni wierny: każdemu odwzorowaniu liniowemu $k^n\to k^m$ odpowiada dokładnie jedna macierz o rozmiarach $m\times n$.
\end{przyklad}


\begin{stwierdzenie}[W pełni wierne funktory odzwierciedlają izomorfizmy]\label{tw:ff-reflect-iso}
Niech $F:\mathcal C\to\mathcal D$ będzie w pełni wierny. Jeśli $f:A\to B$ jest morfizmem w $\mathcal C$ i $F(f)$ jest izomorfizmem w $\mathcal D$, to $f$ jest izomorfizmem w $\mathcal C$.
\end{stwierdzenie}
\begin{proof}
Pełność daje $g:B\to A$ z $F(g)=F(f)^{-1}$. Z funktorialności wynika $F(g\circ f)=F(g)\circ F(f)=\id_{F(A)}=F(\id_A)$, więc wierność daje $g\circ f=\id_A$. Analogicznie $F(f\circ g)=F(f)\circ F(g)=\id_{F(B)}=F(\id_B)$, a wierność daje $f\circ g=\id_B$.
\end{proof}

\begin{przyklad}[Wierność bez pełności]
Funktor zapominania $U:\Grp\to\Set$ jest wierny, lecz niepełny. Na przykład funkcja $q:\mathbb Z\to\mathbb Z$, $q(n)=n^2$, jest funkcją zbiorów, ale nie jest homomorfizmem grup addytywnych, ponieważ $q(1+1)=4\neq 2=q(1)+q(1)$. Pełność wymaga reprezentowania \emph{każdej} funkcji między odpowiednimi zbiorami jako obrazu homomorfizmu, co w tym przykładzie zawodzi.
\end{przyklad}


\subsection{Funktory Hom}

\begin{definicja}[Funktory reprezentowane przez obiekt]\label{def:hom-funktory}
Niech $\mathcal C$ będzie kategorią lokalnie małą i $A\in\Ob(\mathcal C)$. Definiujemy funktor kowariantny $h^A:\mathcal C\to\Set$ przez $h^A(X)=\Hom_{\mathcal C}(A,X)$ oraz $h^A(f)(u)=f\circ u$ dla $f:X\to Y$ i $u:A\to X$. Definiujemy też funktor kontrawariantny $h_A:\mathcal C^{\op}\to\Set$ przez $h_A(X)=\Hom_{\mathcal C}(X,A)$ oraz $h_A(f)(v)=v\circ f$ dla $f:X\to Y$ i $v:Y\to A$.
\end{definicja}

\begin{proof}[Sprawdzenie aksjomatów funktora]
Dla $u:A\to X$ mamy $h^A(\id_X)(u)=\id_X\circ u=u$. Dla $f:X\to Y$, $g:Y\to Z$ zachodzi $h^A(g\circ f)(u)=(g\circ f)\circ u=g\circ(f\circ u)$. Dla $v:Z\to A$ analogicznie $h_A(g\circ f)(v)=v\circ(g\circ f)=(v\circ g)\circ f=h_A(f)(h_A(g)(v))$, a $h_A(\id_X)(w)=w\circ\id_X=w$ dla $w:X\to A$.
\end{proof}

\begin{przyklad}
W $\Set$ funktor $h^1=\Hom_{\Set}(\{*\},-)$ jest naturalnie izomorficzny z funktorem identycznościowym, ponieważ każda funkcja $\{*\}\to X$ jest wyznaczona przez obraz $*$. Natomiast $h^{\{0,1\}}(X)$ jest naturalnie utożsamiany ze zbiorem $X\times X$, ponieważ funkcję na zbiorze dwuelementowym opisują dwa jej obrazy.
\end{przyklad}

\subsection{Zachowywanie i odzwierciedlanie własności morfizmów}

\begin{stwierdzenie}\label{cat:prop:faithful-reflect}
Funktor wierny odzwierciedla monomorfizmy i epimorfizmy. Każdy funktor zachowuje rozszczepione monomorfizmy i rozszczepione epimorfizmy.
\end{stwierdzenie}
\begin{proof}
Niech $F$ będzie wierny i $F(f)$ monomorfizmem. Jeżeli $f\circ u=f\circ v$, to $F(f)\circ F(u)=F(f)\circ F(v)$, skąd $F(u)=F(v)$, a więc $u=v$. Jeżeli $F(f)$ jest epimorfizmem, z $u\circ f=v\circ f$ analogicznie otrzymujemy $F(u)\circ F(f)=F(v)\circ F(f)$, następnie $F(u)=F(v)$ i $u=v$. Jeżeli $r\circ f=\id_A$, to $F(r)\circ F(f)=\id_{F(A)}$; jeżeli $f\circ s=\id_B$, to $F(f)\circ F(s)=\id_{F(B)}$. To daje odpowiednie rozszczepienia w kategorii docelowej.
\end{proof}

\begin{przyklad}[Granice tych własności]\label{cat:ex:functor-limits}
Wierny funktor zapominania $\Top\to\Set$ nie odzwierciedla izomorfizmów: bijekcja ciągła z dwupunktowej przestrzeni dyskretnej do antydyskretnej jest kontrprzykładem z rozdziału~\ref{sec:morfizmy}. Funktor $\Ring\to\Set$ nie zachowuje epimorfizmów (przykład~\ref{przyklad:epi-pierscienie}). Funktor do kategorii z jednym obiektem i jedną strzałką jest pełny tylko wtedy, gdy każda para obiektów źródłowych ma przynajmniej jeden morfizm; dla kategorii z jednym obiektem odpowiadającej nietrywialnej grupie jest zatem pełny, ale niewierny. Pokazuje to niezależność pełności od wierności.
\end{przyklad}

\subsection{Kategorie iloczynowe i bifunktory}

\begin{definicja}[Kategoria iloczynowa i bifunktor]\label{cat:def:bifunctor}
Kategoria $\mathcal C\times\mathcal D$ ma obiekty $(A,B)$ i morfizmy $(f,g):(A,B)\to(A',B')$, gdzie $f:A\to A'$ i $g:B\to B'$. Identyczności i składanie definiujemy współrzędnie:
\[
(f',g')\circ(f,g)=(f'\circ f,g'\circ g),
\qquad \id_{(A,B)}=(\id_A,\id_B).
\]
Aksjomaty kategorii wynikają z odpowiednich równości w każdej współrzędnej. Bifunktor to funktor $F:\mathcal C\times\mathcal D\to\mathcal E$.
\end{definicja}

\begin{stwierdzenie}[Bifunktor Hom]\label{cat:prop:hom-bifunctor}
Dla kategorii lokalnie małej $\mathcal C$ istnieje bifunktor
\[
\Hom_{\mathcal C}:\mathcal C^{\op}\times\mathcal C\longrightarrow\Set.
\]
Na obiektach przyjmuje wartości $\Hom_{\mathcal C}(A,B)$, a dla $a:A'\to A$, $b:B\to B'$ i $u:A\to B$ działa przez $u\mapsto b\circ u\circ a$.
\end{stwierdzenie}
\begin{proof}
Obraz ma poprawny typ $A'\to B'$. Identyczności po obu stronach nie zmieniają $u$. Dla kolejnych strzałek $a':A''\to A'$ i $b':B'\to B''$ dwukrotne zastosowanie bifunktora daje $b'\circ(b\circ u\circ a)\circ a'=(b'\circ b)\circ u\circ(a\circ a')$, czyli dokładnie obraz złożenia w kategorii iloczynowej. Ustalając pierwszy lub drugi argument, odzyskujemy funktory z definicji~\ref{def:hom-funktory}.
\end{proof}

\begin{przyklad}[Iloczyn tensorowy jako bifunktor]\label{cat:ex:tensor-bifunctor}
Niech $R$ będzie pierścieniem przemiennym z jedynką. Korzystamy z własności uniwersalnej iloczynu tensorowego z rozdziału 18 skryptu \emph{Algebra abstrakcyjna}. Dla homomorfizmów $f:M\to M'$ i $g:N\to N'$ odwzorowanie $(x,y)\mapsto f(x)\otimes g(y)$ jest dwuliniowe, więc indukuje jedyny homomorfizm $f\otimes g:M\otimes_R N\to M'\otimes_R N'$. Dla kolejnych homomorfizmów $f',g'$ mamy
\[
((f'\otimes g')\circ(f\otimes g))(x\otimes y)
=(f'\circ f)(x)\otimes(g'\circ g)(y).
\]
Tensory proste generują iloczyn tensorowy jako moduł, zatem zachowana jest kompozycja; identyczności sprawdzamy tak samo na tensorach prostych. Otrzymujemy bifunktor kowariantny.

Funktor $M\mapsto M\otimes_{\mathbb Z}\mathbb Z/2\mathbb Z$ nie zachowuje monomorfizmów: iniekcję $2:\mathbb Z\to\mathbb Z$ przeprowadza na odwzorowanie zerowe niezerowego modułu $\mathbb Z/2\mathbb Z$ w siebie. Wykorzystujemy tutaj izomorfizm $\mathbb Z\otimes_{\mathbb Z}\mathbb Z/2\mathbb Z\cong\mathbb Z/2\mathbb Z$ oraz relację $2\otimes y=1\otimes 2y=0$.
\end{przyklad}

\subsection{Funktory w geometrii i topologii}

\begin{przyklad}[Spektrum jako funktor kontrawariantny]\label{cat:ex:spec-functor}
Dla pierścienia przemiennego $R$ z jedynką rozważmy przestrzeń $\operatorname{Spec} R$ ideałów pierwszych z topologią Zariskiego (rozdział 4 skryptu \emph{Geometria algebraiczna}). Homomorfizm $\varphi:R\to S$ zachowujący jedynkę wyznacza funkcję $\operatorname{Spec} S\to\operatorname{Spec} R$, $\mathfrak p\mapsto\varphi^{-1}(\mathfrak p)$. Przeciwobraz ideału pierwszego jest ideałem pierwszym: nie zawiera jedynki, a z $\varphi(ab)=\varphi(a)\varphi(b)\in\mathfrak p$ wynika $\varphi(a)\in\mathfrak p$ lub $\varphi(b)\in\mathfrak p$. Funkcja jest ciągła, gdyż przeciwobraz otwartego zbioru bazowego $D_R(r)=\{\mathfrak q:r\notin\mathfrak q\}$ to $D_S(\varphi(r))$. Dla $\psi:S\to T$ i $\mathfrak p\in\operatorname{Spec} T$ mamy
\[
(\psi\circ\varphi)^{-1}(\mathfrak p)=\varphi^{-1}(\psi^{-1}(\mathfrak p)).
\]
Identyczność działa przez identyczność. Otrzymujemy $\operatorname{Spec}:\Ring^{\op}\to\Top$. W tym miejscu rozpatrujemy jedynie przestrzeń topologiczną; konstrukcję morfizmu snopów strukturalnych i funktora do kategorii schematów afinicznych pozostawiamy do dalszej części.
\end{przyklad}

\begin{przyklad}[Homologia kompleksów]\label{cat:ex:homology-functor}
Dla kompleksu łańcuchowego $C_\bullet$ modułów nad ustalonym pierścieniem oznaczmy przez $H_n(C_\bullet)=\ker d_n/\operatorname{im}d_{n+1}$ jego homologię stopnia $n$. Morfizm kompleksów $f:C_\bullet\to D_\bullet$ spełnia $d_n^D f_n=f_{n-1}d_n^C$. Przeprowadza więc cykle na cykle, a brzegi na brzegi: jeśli $x=d_{n+1}^C y$, to $f_n(x)=d_{n+1}^D f_{n+1}(y)$. Wzór $H_n(f)([x])=[f_n(x)]$ jest dobrze określony na klasach ilorazowych. Zachowanie identyczności i składania wynika z $\id([x])=[x]$ oraz $H_n(g\circ f)([x])=[g_n(f_n(x))]$. Jest to konstrukcja z rozdziału 3 skryptu \emph{Topologia algebraiczna}, przedstawiona tutaj jako funktor kowariantny.
\end{przyklad}

\begin{przyklad}[Cofnięcie form różniczkowych i kohomologia de Rhama]\label{cat:ex:derham-functor}
Dla gładkiego $f:M\to N$ cofnięcie $f^*:\Omega^q(N)\to\Omega^q(M)$ określa wzór
\[
(f^*\omega)_x(v_1,\ldots,v_q)=
\omega_{f(x)}(\mathrm df_x(v_1),\ldots,\mathrm df_x(v_q)).
\]
Dla $g:N\to P$ reguła łańcuchowa $\mathrm d(g\circ f)_x=\mathrm dg_{f(x)}\circ\mathrm df_x$ daje $(g\circ f)^*=f^*\circ g^*$ przez podstawienie do powyższego wzoru. Cofnięcie identycznością pozostawia formę bez zmian. Korzystamy ponadto z udowodnionej w rozdziale 10 skryptu \emph{Geometria i topologia różniczkowa} zgodności $d(f^*\omega)=f^*(d\omega)$ z różniczką zewnętrzną. Wobec tego cofnięcie zachowuje formy zamknięte i dokładne oraz indukuje $H^q_{\mathrm{dR}}(f):H^q_{\mathrm{dR}}(N)\to H^q_{\mathrm{dR}}(M)$, $[\omega]\mapsto[f^*\omega]$. Równości dla cofnięć dają te same równości na klasach kohomologii, więc zarówno formy ustalonego stopnia, jak i kohomologia de Rhama wyznaczają funktory kontrawariantne do kategorii rzeczywistych przestrzeni wektorowych.
\end{przyklad}

